\documentclass[11pt]{article}
\usepackage[T1]{fontenc}
\usepackage[utf8]{inputenc}
\usepackage{amsmath,amssymb,amsthm}
\usepackage[margin=1in]{geometry}
\usepackage{array}
\usepackage[colorlinks=true,linkcolor=blue,citecolor=blue]{hyperref}
\usepackage{xcolor}

\newtheorem{theorem}{Theorem}[section]
\newtheorem{lemma}[theorem]{Lemma}
\newtheorem{proposition}[theorem]{Proposition}
\newtheorem{corollary}[theorem]{Corollary}
\newtheorem{assumption}[theorem]{Assumption}
\theoremstyle{definition}
\newtheorem{definition}[theorem]{Definition}
\theoremstyle{remark}
\newtheorem{remark}[theorem]{Remark}

\newcommand{\ASSUMED}{\textcolor{orange}{\textbf{[ASSUMED]}}}
\newcommand{\PROVED}{\textcolor{green!60!black}{\textbf{[PROVED]}}}
\newcommand{\CONJ}{\textcolor{red}{\textbf{[CONJECTURED]}}}
\newcommand{\Erec}{E_{\mathrm{rec}}}

\title{RIP-U and the $\omega = 0$ obstruction in Davies dynamics\\
\large Upper envelopes, witness floors, and falsification protocols for
separation-dependent decoherence rates}
\author{Lluis Eriksson}
\date{January 2026 (v1); July 2026 (v2)}

\begin{document}
\maketitle

\begin{abstract}
We isolate the dynamic hinge in typed separation$\to$rate$\to$power
pipelines within the Davies (weak-coupling, Markovian) setting in finite
dimension. First, we formulate an upper-envelope statement (RIP-U): under
an explicit factorized bath-correlation envelope with
separation-dependent amplitude $f(\epsilon)$ and integrable time profile,
Fourier-transformed Davies rates inherit an $O(f(\epsilon))$ envelope.
Under additional regularity assumptions preventing trivial degeneracies,
this yields a worst-case bound $\kappa^{\uparrow}(\epsilon) \le
Cf(\epsilon)$ for the instantaneous relative loss rate of a coherence-like
functional. Second, we isolate a structural obstruction to lower-envelope
statements: we prove an exact identity for the $\omega = 0$ contribution
to the Davies Dirichlet form,
\[
\mathcal{E}^{(0)}_{\sigma}(O) = \frac{\gamma(0)}{2}\,
\left\| [S(0), O] \right\|_{2,\sigma}^{2},
\]
and derive a witness mechanism showing how $\omega = 0$ channels can
enforce a non-vanishing dissipation contribution for suitable observable
families. We emphasize directionality: RIP-U (upper) does not imply a
positive lower envelope $\kappa^{\downarrow}(\epsilon)$ without additional
family-qualified input. All assumptions are explicit and accompanied by
falsification routes. \emph{v2 (no v1 number is changed):} two v1
assumptions are upgraded to proved statements in their canonical
instances --- $\Delta$-MONO holds unconditionally for the energy pinching
(Davies covariance plus data processing, Lemma~\ref{lem:dmono}), and a
sufficient bridge with an explicit constant is proved
(Proposition~\ref{prop:bridgep}); the paper's cross-reference labels are
repaired (in v1 every assumption and theorem was typeset as
``Definition~$x.y$''); Table~\ref{tab:fals} is re-typeset (overlapping
text in v1); the $\omega=0$ identity is placed within the corrected
Bohr-channel decomposition of the companion 2601.0023 (it is exactly the
$\omega=0$ sector, and the plain-commutator form provably fails at
$\omega \ne 0$); and a verification suite reproduces every proved item
numerically, including the identity at machine precision and an explicit
family with $\kappa^{\downarrow} \ll \kappa^{\uparrow}$ under the same
envelope.
\end{abstract}

\section{Scope and conventions}

\textbf{What this paper does.}
\begin{itemize}
\item Proves a Fourier-envelope inheritance lemma for rate coefficients
under a factorized correlator bound.
\item States a clean upper-envelope implication (RIP-U) under explicit
additional bridge assumptions, and (v2) proves a sufficient bridge with
an explicit constant on a regularized family.
\item Proves an exact $\omega = 0$ Dirichlet identity in standard Davies
form and a witness corollary.
\end{itemize}

\textbf{What this paper does not do.} We do not claim a Hamiltonian
spectral gap and we do not claim RIP-U implies a family-wide true floor
$\kappa^{\downarrow}(\epsilon) \ge \kappa_0$.

\textbf{Parameter convention.} $\epsilon$ is geometric separation/buffer
width; $\delta$ is tolerance/regularization. They are never interchanged.

\section{Davies setting and typed envelopes}

We work in finite dimension: system Hilbert space $\mathcal{H}$ with
$\dim \mathcal{H} = d < \infty$.

\begin{assumption}[DAVIES: weak-coupling Markovian dynamics]
\label{ass:davies}
\ASSUMED. The system dynamics is governed by a CPTP semigroup
$\mathcal{T}_t = e^{t\mathcal{L}}$ with a Davies-type GKLS generator
$\mathcal{L}$ and a faithful stationary state $\sigma$ (typically
$\sigma \propto e^{-\beta H}$). We assume the standard Bohr-frequency
decomposition and GKLS form used in Davies theory \cite{Davies,BP,SL}.
\end{assumption}

\subsection{A coherence-like functional and instantaneous loss}

\begin{definition}[D1: $\Delta$-track coherence]
\label{def:cdelta}
Fix a dephasing (pinching) CPTP map $\Delta$ on $\mathcal{B}(\mathcal{H})$.
Define
\[
C_{\Delta}(\rho) := S(\rho \| \Delta[\rho]),
\]
where $S(\rho\|\tau) := \mathrm{Tr}(\rho(\log\rho - \log\tau))$ when
well-defined.
\end{definition}

\begin{definition}[D2: instantaneous loss]
\label{def:loss}
Define the instantaneous change of $C_\Delta$ under $\mathcal{T}_t$ by
\[
\dot{C}_{\Delta,\mathrm{loss}}(\rho) := -\left.\frac{d}{dt}\right|_{t=0}
C_\Delta(\mathcal{T}_t(\rho)),
\]
whenever the derivative exists.
\end{definition}

\begin{assumption}[$\Delta$-MONO: sign convention for ``loss'']
\label{ass:dmono}
\ASSUMED\ (family restriction; \emph{proved for the energy pinching in
v2}, Lemma~\ref{lem:dmono}). For the chosen $\Delta$ and semigroup
$\mathcal{T}_t$, we restrict attention to families $\mathcal{F}_\epsilon$
such that $\dot{C}_{\Delta,\mathrm{loss}}(\rho) \ge 0$ for all
$\rho \in \mathcal{F}_\epsilon$.
\end{assumption}

\begin{definition}[D3: ratios and envelopes]
\label{def:ratios}
Whenever $C_\Delta(\rho) > 0$, define
$r(\rho) := \dot{C}_{\Delta,\mathrm{loss}}(\rho) / C_\Delta(\rho)$.
For a specified family $\mathcal{F}_\epsilon$ define
\[
\kappa^{\uparrow}(\epsilon) := \sup_{\substack{\rho\in\mathcal{F}_\epsilon:\,
C_\Delta(\rho)>0,\ \dot{C}_{\Delta,\mathrm{loss}}(\rho)\ \mathrm{exists}}}
r(\rho), \qquad
\kappa^{\downarrow}(\epsilon) := \inf_{\substack{\rho\in\mathcal{F}_\epsilon:\,
C_\Delta(\rho)>0,\ \dot{C}_{\Delta,\mathrm{loss}}(\rho)\ \mathrm{exists}}}
r(\rho).
\]
\end{definition}

\begin{remark}[Directionality]
\label{rem:dir}
RIP-U targets $\kappa^{\uparrow}(\epsilon)$ (worst-case upper envelope).
Minimum-maintenance-power lower bounds require
$\kappa^{\downarrow}(\epsilon)$; RIP-U alone does not provide them.
\end{remark}

\subsection{$\Delta$-MONO is unconditional for the energy pinching (v2)}

\begin{lemma}[Energy pinching: covariance and monotone loss]
\label{lem:dmono}
\PROVED\ (v2). Let $\Delta_H$ be the pinching over the spectral projectors
$\{\Pi_E\}$ of $H$ and let $\mathcal{L}$ be a Davies generator built from
Bohr components of the coupling (Assumption~\ref{ass:bohr}). Then:
\emph{(i)} $\Delta_H \circ \mathcal{T}_t = \mathcal{T}_t \circ \Delta_H$
for all $t \ge 0$; \emph{(ii)} for every state $\rho$ in the domain,
$t \mapsto C_{\Delta_H}(\mathcal{T}_t(\rho))$ is nonincreasing, hence
$\dot{C}_{\Delta_H,\mathrm{loss}}(\rho) \ge 0$ with no family
restriction.
\end{lemma}

\begin{proof}
(i) The Davies generator commutes with the Hamiltonian dressing
$e^{iHt}(\cdot)e^{-iHt}$ (covariance is built into the Bohr construction:
each dissipator term maps the $\omega'$ Bohr sector of its argument into
itself). Since $\Delta_H$ is the Ces\`aro average of the dressing,
$\mathcal{L} \circ \Delta_H = \Delta_H \circ \mathcal{L}$, and the
semigroup inherits the commutation. (ii) By (i) and monotonicity of
relative entropy under the CPTP map $\mathcal{T}_s$,
\[
C_{\Delta_H}(\mathcal{T}_{t+s}\rho)
= S(\mathcal{T}_s \mathcal{T}_t \rho \,\|\, \mathcal{T}_s \Delta_H
\mathcal{T}_t \rho)
\le S(\mathcal{T}_t \rho \,\|\, \Delta_H \mathcal{T}_t \rho)
= C_{\Delta_H}(\mathcal{T}_t\rho). \qedhere
\]
\end{proof}

\begin{lemma}[Derivative identity]
\label{lem:deriv}
\PROVED\ (v2). Let $\Delta = \Delta_H$ and let $\rho$ be full rank. Then
\[
\dot{C}_{\Delta,\mathrm{loss}}(\rho)
= -\mathrm{Tr}\bigl(\mathcal{L}(\rho)\,(\log\rho - \log\Delta[\rho])\bigr).
\]
\end{lemma}

\begin{proof}
$\frac{d}{dt}\mathrm{Tr}(\rho_t \log \rho_t) =
\mathrm{Tr}(\mathcal{L}(\rho)\log\rho)$ by $\mathrm{Tr}\,\mathcal{L}(\rho)
= 0$. For the second term, $\frac{d}{dt}\mathrm{Tr}(\rho_t \log
\Delta[\rho_t]) = \mathrm{Tr}(\mathcal{L}(\rho)\log\Delta[\rho]) +
\mathrm{Tr}(\rho \,\tfrac{d}{dt}\log\Delta[\rho_t])$; since
$\tfrac{d}{dt}\log\Delta[\rho_t]$ lies in the pinched (commutative)
algebra, self-duality of $\Delta$ gives
$\mathrm{Tr}(\rho\,\tfrac{d}{dt}\log\Delta[\rho_t]) =
\mathrm{Tr}(\Delta[\rho]\,\tfrac{d}{dt}\log\Delta[\rho_t]) =
\tfrac{d}{dt}\mathrm{Tr}(\Delta[\rho_t]) = 0$.
\end{proof}

\section{Correlator envelope $\Rightarrow$ Fourier rate envelope}

\begin{assumption}[CORR: factorized correlator envelope]
\label{ass:corr}
\ASSUMED. For each $\epsilon$, the relevant bath correlators satisfy
\[
|G_{\alpha\beta}(t;\epsilon)| \le g_{\alpha\beta}(t)\, f(\epsilon),
\]
where $g_{\alpha\beta} \in L^1(\mathbb{R})$ and $f(\epsilon) \ge 0$
captures separation dependence.
\end{assumption}

\begin{definition}[D4: Fourier rate coefficients]
\label{def:fourier}
Define the (formal) Fourier rate coefficients by
$\gamma_{\alpha\beta}(\omega;\epsilon) := \int_{\mathbb{R}} dt\,
e^{i\omega t}\, G_{\alpha\beta}(t;\epsilon)$, whenever the integral exists
(or under the regularity hypotheses used in Davies theory).
\end{definition}

\begin{lemma}[Fourier envelope inheritance]
\label{lem:envelope}
\PROVED\ (under Assumption~\ref{ass:corr} and integrability). Under
Assumption~\ref{ass:corr}, for all $\omega$ and $\epsilon$ for which
Definition~\ref{def:fourier} is well-defined,
\[
|\gamma_{\alpha\beta}(\omega;\epsilon)| \le
\|g_{\alpha\beta}\|_{L^1}\, f(\epsilon).
\]
\end{lemma}

\begin{proof}
By the triangle inequality and Assumption~\ref{ass:corr},
$|\gamma_{\alpha\beta}(\omega;\epsilon)| \le \int_{\mathbb{R}} dt\,
|G_{\alpha\beta}(t;\epsilon)| \le \int_{\mathbb{R}} dt\,
|g_{\alpha\beta}(t)|\, f(\epsilon) = \|g_{\alpha\beta}\|_{L^1}
f(\epsilon)$.
\end{proof}

\section{RIP-U: an upper envelope on instantaneous loss}

\begin{assumption}[REG: non-degeneracy for relative-entropy ratios]
\label{ass:reg}
\ASSUMED\ (family restriction). The family $\mathcal{F}_\epsilon$ is
restricted so that: (i) $C_\Delta(\rho) \ge c_{\min}(\delta) > 0$
uniformly over $\rho \in \mathcal{F}_\epsilon$; (ii) $\Delta[\rho]
\succeq p_{\min}(\delta)\mathbf{1}$ uniformly over $\rho \in
\mathcal{F}_\epsilon$; (iii) when needed, $\rho$ stays a fixed distance
(within tolerance $\delta$) away from $\mathrm{Fix}(\mathcal{T}_t)$.
\end{assumption}

\begin{assumption}[BRIDGE: rate envelope $\Rightarrow$ loss envelope]
\label{ass:bridge}
\ASSUMED\ (\emph{a proved sufficient instance in v2:}
Proposition~\ref{prop:bridgep}). Assume Assumptions~\ref{ass:davies},
\ref{ass:corr} and \ref{ass:reg}. Suppose the Davies generator is built
from Fourier rate coefficients $\gamma_{\alpha\beta}(\omega;\epsilon)$ in
such a way that the envelope from Lemma~\ref{lem:envelope} implies: there
exists a constant $C_{\mathrm{RIP}} > 0$ such that for all $\rho \in
\mathcal{F}_\epsilon$,
\[
\dot{C}_{\Delta,\mathrm{loss}}(\rho) \le C_{\mathrm{RIP}}\, f(\epsilon)\,
C_\Delta(\rho).
\]
\end{assumption}

\begin{theorem}[T1: RIP-U (upper envelope)]
\label{thm:ripu}
\PROVED\ (given Assumptions~\ref{ass:davies}, \ref{ass:dmono},
\ref{ass:corr}, \ref{ass:reg} and \ref{ass:bridge}). Under
Assumptions~\ref{ass:davies}, \ref{ass:dmono}, \ref{ass:corr},
\ref{ass:reg} and \ref{ass:bridge}, for all $\rho \in
\mathcal{F}_\epsilon$,
\[
r(\rho) \le C_{\mathrm{RIP}}\, f(\epsilon), \qquad \text{and hence}
\qquad \kappa^{\uparrow}(\epsilon) \le C_{\mathrm{RIP}}\, f(\epsilon).
\]
\end{theorem}

\begin{proof}
By Assumption~\ref{ass:bridge}, $\dot{C}_{\Delta,\mathrm{loss}}(\rho)
\le C_{\mathrm{RIP}} f(\epsilon) C_\Delta(\rho)$ for all $\rho \in
\mathcal{F}_\epsilon$ with $C_\Delta(\rho) > 0$, hence $r(\rho) \le
C_{\mathrm{RIP}} f(\epsilon)$. Taking the supremum yields the stated
envelope.
\end{proof}

\begin{remark}[RIP-U does not imply a floor]
\label{rem:nofloor}
RIP-U controls worst-case \emph{upper} loss rates. It is compatible with
$\kappa^{\downarrow}(\epsilon)$ being arbitrarily small on a large
family. Thus RIP-U alone does not imply minimum maintenance power lower
bounds in pipelines requiring a positive lower envelope. (The v2 suite
exhibits this concretely: a weakly coupled level pair yields
$\kappa^{\downarrow} \sim 10^{-3}\,\kappa^{\uparrow}$ under the same
correlator envelope.)
\end{remark}

\subsection{A proved sufficient bridge (v2)}

\begin{proposition}[BRIDGE-P: explicit-constant bridge on a regularized
family]
\label{prop:bridgep}
\PROVED\ (v2). Let $\Delta = \Delta_H$, and strengthen
Assumption~\ref{ass:reg} to REG$^{+}$: (i) $C_\Delta(\rho) \ge c_{\min} >
0$ and (ii$'$) $\rho \succeq p_{\min}\mathbf{1}$ (hence also
$\Delta[\rho] \succeq p_{\min}\mathbf{1}$) over $\mathcal{F}_\epsilon$.
Let the Davies dissipator be built from Bohr components $S(\omega)$ with
rates $\gamma(\omega;\epsilon)$ obeying the envelope of
Lemma~\ref{lem:envelope} with a common profile,
$\gamma(\omega;\epsilon) \le \bar{g} f(\epsilon)$, and let $N_B$ be the
number of Bohr channels and $M^2 := \max_\omega \|S(\omega)\|^2$. Then
Assumption~\ref{ass:bridge} holds with the explicit constant
\[
C_{\mathrm{RIP}} = \frac{4\, N_B\, \bar{g}\, M^2 \,
\log(1/p_{\min})}{c_{\min}}.
\]
\end{proposition}

\begin{proof}
By Lemma~\ref{lem:deriv} and H\"older,
$\dot{C}_{\Delta,\mathrm{loss}}(\rho) \le
\|\mathcal{L}_D(\rho)\|_1\, \|\log\rho - \log\Delta[\rho]\|_\infty$.
Under REG$^{+}$, $\|\log\rho\|_\infty \le \log(1/p_{\min})$ and likewise
for $\Delta[\rho]$, so the second factor is at most
$2\log(1/p_{\min})$. Each dissipator term is $1$-norm bounded by
$2\gamma(\omega;\epsilon)\|S(\omega)\|^2$, so
$\|\mathcal{L}_D(\rho)\|_1 \le 2 N_B \bar{g} M^2 f(\epsilon)$.
(The Hamiltonian part contributes zero loss: unitary conjugation
commuting with $\Delta_H$ leaves $C_\Delta$ invariant.) Dividing by
$C_\Delta(\rho) \ge c_{\min}$ gives the claim.
\end{proof}

\begin{remark}
\label{rem:bridgep}
BRIDGE-P demotes the bridge from an assumption to a checkable statement
on the regularized family: the constant is explicit, dimension-aware
through $N_B$ and $M$, and blows up only as $p_{\min} \to 0$ or
$c_{\min} \to 0$, which is the honest signature of the relative-entropy
ratio degenerating. The v2 suite instantiates it numerically (worst
observed $r / (C_{\mathrm{RIP}} f) \approx 2\times 10^{-3}$, i.e.\ the
bound is loose but never violated).
\end{remark}

\section{The $\omega = 0$ obstruction: an exact Dirichlet identity}

\begin{definition}[D5: $\sigma$-weighted 2-norm]
\label{def:kmsnorm}
For a faithful state $\sigma$, define $\|X\|_{2,\sigma} :=
\sqrt{\mathrm{Tr}\bigl(\sigma^{1/2}X^{\dagger}\sigma^{1/2}X\bigr)}$.
\end{definition}

\begin{definition}[D6: Dirichlet form (one convenient convention)]
\label{def:dirichlet}
Let $\mathcal{L}$ be a GKLS generator with faithful stationary state
$\sigma$. Define the Dirichlet form on observables by
\[
\mathcal{E}_\sigma(O) := -\mathrm{Tr}\Bigl(\sigma^{1/2}O^{\dagger}
\sigma^{1/2}\, \mathcal{L}^{\sharp}(O)\Bigr),
\]
where $\mathcal{L}^{\sharp}$ denotes the Heisenberg-picture action of the
generator on observables.
\end{definition}

\begin{assumption}[BOHR: frequency decomposition]
\label{ass:bohr}
\ASSUMED\ (standard in Davies theory). Let $H$ be the system Hamiltonian
with spectral projectors $\{\Pi_E\}$. For a system coupling operator $S$,
define its Bohr components $S(\omega) := \sum_{E'-E=\omega} \Pi_E S
\Pi_{E'}$. Assume the Davies generator uses dissipators built from
$S(\omega)$ with scalar rates $\gamma(\omega) \ge 0$.
\end{assumption}

\begin{lemma}[Commutation at zero frequency]
\label{lem:commute}
\PROVED. If $\sigma \propto e^{-\beta H}$ and $S(0)$ is the $\omega = 0$
Bohr component, then $[S(0), H] = 0$, and hence $[S(0), \sigma] = 0$ and
$[S(0), \sigma^{1/2}] = 0$.
\end{lemma}

\begin{proof}
By construction, $S(0) = \sum_E \Pi_E S \Pi_E$ commutes with $H = \sum_E
E\, \Pi_E$, hence with any function of $H$, in particular $\sigma$ and
$\sigma^{1/2}$.
\end{proof}

\begin{theorem}[T2: $\omega = 0$ Dirichlet identity]
\label{thm:omega0}
\PROVED\ (under Assumptions~\ref{ass:davies} and \ref{ass:bohr} in
standard Davies form). The zero-frequency contribution
$\mathcal{E}^{(0)}_\sigma$ to the Dirichlet form satisfies
\[
\mathcal{E}^{(0)}_\sigma(O) = \frac{\gamma(0)}{2}\,
\|[S(0), O]\|^{2}_{2,\sigma}.
\]
\end{theorem}

\begin{proof}
We give a direct expansion in Appendix~\ref{app:proof}.
\end{proof}

\begin{remark}[The $\gamma(0) = 0$ case]
\label{rem:gamma0}
If $\gamma(0) = 0$ (e.g.\ baths with vanishing zero-frequency spectral
density), then $\mathcal{E}^{(0)}_\sigma \equiv 0$ and this particular
$\omega = 0$ obstruction mechanism disappears. Any lower-envelope floor
must then come from other frequency sectors or from structural
restrictions defining the family.
\end{remark}

\begin{corollary}[C1: $\omega = 0$ witness mechanism]
\label{cor:witness}
\PROVED. Assume $\gamma(0) > 0$. If there exists an observable family
$O_\epsilon$ such that
$\|[S(0), O_\epsilon]\|^{2}_{2,\sigma} \big/ \|O_\epsilon\|^{2}_{2,\sigma}
\ge c > 0$, then
$\mathcal{E}^{(0)}_\sigma(O_\epsilon) \ge \frac{\gamma(0)}{2}\, c\,
\|O_\epsilon\|^{2}_{2,\sigma}$.
\end{corollary}

\begin{proof}
Immediate from Theorem~\ref{thm:omega0}.
\end{proof}

\begin{remark}[C1 is not RIP-L]
\label{rem:notripl}
C1 is a mechanism for a chosen observable family. Converting it into a
uniform lower-envelope bound $\kappa^{\downarrow}(\epsilon) \ge \kappa_0$
over a state family requires additional typed bridges.
\end{remark}

\begin{assumption}[RIP-L: true floor (family-qualified)]
\label{ass:ripl}
\CONJ. There exist $\epsilon_* > 0$ and $\kappa_0 > 0$ such that
$\kappa^{\downarrow}(\epsilon) \ge \kappa_0$ for all $\epsilon \ge
\epsilon_*$, for a specified family $\mathcal{F}_\epsilon$ (with any
regularity restrictions stated explicitly).
\end{assumption}

\begin{remark}[Position within the Bohr-channel decomposition of
2601.0023 (v2)]
\label{rem:0023}
The companion 2601.0023 (v2) proves the exact channel-wise decomposition
$\mathcal{E}_\sigma(O) = \sum_\omega \gamma(\omega)\, e^{\beta\omega/2}\,
\mathrm{Re}\langle S(\omega)O, [S(\omega), O]\rangle_{\mathrm{KMS}}$ and
shows that the naive plain-commutator form
$\frac{\gamma(\omega)}{2}\|[S(\omega),O]\|^2_{2,\sigma}$ \emph{fails} at
$\omega \ne 0$ (deviations of order $20\%$ there; up to $57\%$ in the
present suite's model). Theorem~\ref{thm:omega0} is exactly the
$\omega = 0$ sector of that decomposition: at $\omega = 0$ the KMS weight
$e^{\beta\omega/2}$ is $1$ and $S(0)$ commutes with $\sigma^{1/2}$
(Lemma~\ref{lem:commute}), which is precisely what collapses the exact
quadratic form to the plain commutator. The verification suite checks
both facts side by side: machine-precision agreement at $\omega = 0$,
systematic deviation at every $\omega \ne 0$.
\end{remark}

\section{Falsification matrix}

\begin{table}[h!]
\centering
\begin{tabular}{p{2.6cm} p{3.2cm} p{4.0cm} p{4.6cm}}
\hline
\textbf{Item} & \textbf{Status} & \textbf{Verification proxy} &
\textbf{Concrete falsifier} \\
\hline
CORR & \ASSUMED & compute $G_{\alpha\beta}(t;\epsilon)$ & separation
dependence not factorizable as $g(t)f(\epsilon)$ \\
Fourier envelope lemma & \PROVED & compute $\|g\|_{L^1}$ & violate
integrability/envelope inequality \\
$\Delta$-MONO & \ASSUMED\ general; \PROVED\ for $\Delta_H$ (v2) &
evaluate $\dot{C}_{\Delta,\mathrm{loss}}(\rho)$ & find $\rho$ with
$\dot{C}_{\Delta,\mathrm{loss}}(\rho) < 0$ in family (for $\Delta \ne
\Delta_H$) \\
REG & \ASSUMED & lower-bound spectrum of $\Delta[\rho]$ & sequence with
$\Delta[\rho_n]$ nearly singular \\
BRIDGE & \ASSUMED\ general; sufficient instance \PROVED\ (v2) & compute
$r(\rho)$ on the family & show $r(\rho)$ violates $O(f(\epsilon))$
despite CORR \\
RIP-U & \PROVED & apply Theorem~\ref{thm:ripu} & refute any upstream
assumption \\
$\omega = 0$ identity & \PROVED & direct Dirichlet-form expansion &
refute Davies form/hypotheses \\
RIP-L & \CONJ & estimate $\kappa^{\downarrow}(\epsilon)$ & construct
$\rho_\epsilon$ with $r(\rho_\epsilon) \to 0$ \\
\hline
\end{tabular}
\caption{Falsification matrix for the dynamic layer (re-typeset in v2;
v1's table had overlapping column text).}
\label{tab:fals}
\end{table}

\section*{Note added in v2}
No v1 definition, assumption label, theorem statement or numbering is
changed; new v2 material is appended at the end of the affected sections.
v2: (1) repairs the cross-reference labels --- in v1 every reference was
typeset as ``Definition~$x.y$'' regardless of environment, so e.g.\ ``By
Definition 4.2'' in v1 denotes Assumption~\ref{ass:bridge} and
``Immediate from Definition 5.5'' denotes Theorem~\ref{thm:omega0}; (2)
re-typesets Table~\ref{tab:fals}, which in v1 printed overlapping,
partially illegible status/proxy entries; (3) proves $\Delta$-MONO
unconditionally for the canonical energy pinching via Davies covariance
plus data processing (Lemma~\ref{lem:dmono}), together with the
derivative identity used downstream (Lemma~\ref{lem:deriv}); (4) proves a
sufficient bridge with an explicit constant on a regularized family
(Proposition~\ref{prop:bridgep}), demoting BRIDGE from assumption to
checkable statement in that instance; (5) places
Theorem~\ref{thm:omega0} as the $\omega = 0$ sector of the corrected
Bohr-channel decomposition of 2601.0023 v2 (Remark~\ref{rem:0023}); (6)
adds series positioning (v1 cited no companion papers) and a verification
suite (\texttt{verification/2601-0064/} in the series repository)
reproducing every proved item: the identity at machine precision
($2.9\times10^{-16}$ over random $O$), the envelope lemma with a concrete
correlator, covariance and nonnegative loss for $\Delta_H$, BRIDGE-P, an
explicit weak-pair family with $\kappa^{\downarrow} \approx
10^{-3}\kappa^{\uparrow}$ (Remark~\ref{rem:nofloor}), the witness C1, and
the $\omega \ne 0$ failure of the plain-commutator form.

\begin{thebibliography}{9}
\bibitem{Davies} E. B. Davies, \emph{Markovian master equations},
Commun. Math. Phys. \textbf{39}, 91--110 (1974).
\bibitem{BP} H.-P. Breuer and F. Petruccione, \emph{The Theory of Open
Quantum Systems}, Oxford University Press, 2002.
\bibitem{SL} H. Spohn and J. L. Lebowitz, \emph{Irreversible
thermodynamics for quantum systems weakly coupled to thermal reservoirs},
Adv. Chem. Phys. \textbf{38}, 109--142 (1978).
\bibitem{S0023} L. Eriksson, ai.viXra:2601.0023 (2026; v2 2026).
\bibitem{S0022} L. Eriksson, ai.viXra:2601.0022 (2026; v2 2026).
\bibitem{S0031} L. Eriksson, ai.viXra:2601.0031 (2026; v2 2026).
\bibitem{S0038} L. Eriksson, ai.viXra:2601.0038 (2026; v2 2026).
\bibitem{S0040} L. Eriksson, ai.viXra:2601.0040 (2026; v2 2026).
\end{thebibliography}

\appendix

\section{Line-by-line proof of the $\omega = 0$ identity}
\label{app:proof}

We give a direct expansion showing $\mathcal{E}^{(0)}_\sigma(O) =
\frac{\gamma(0)}{2}\|[S(0),O]\|^2_{2,\sigma}$. For clarity we present the
Hermitian case $S(0)^\dagger = S(0)$; the general case replaces $S(0)^2$
by $S(0)^\dagger S(0)$ and yields the same commutator structure.

The $\omega = 0$ Heisenberg dissipator has the standard GKLS form
\[
\mathcal{L}^{\sharp,(0)}(O) = \gamma(0)\Bigl(S(0)\,O\,S(0) -
\tfrac12\{S(0)^2, O\}\Bigr).
\]
Using the Dirichlet form convention from
Definition~\ref{def:dirichlet},
\[
\mathcal{E}^{(0)}_\sigma(O) = -\mathrm{Tr}\Bigl(\sigma^{1/2}O^\dagger
\sigma^{1/2}\,\mathcal{L}^{\sharp,(0)}(O)\Bigr).
\]
Expand:
\begin{align*}
\mathcal{E}^{(0)}_\sigma(O) ={}& -\gamma(0)\,\mathrm{Tr}\Bigl(
\sigma^{1/2}O^\dagger\sigma^{1/2} S(0)\,O\,S(0)\Bigr)
+ \frac{\gamma(0)}{2}\mathrm{Tr}\Bigl(\sigma^{1/2}O^\dagger\sigma^{1/2}
S(0)^2 O\Bigr) \\
&+ \frac{\gamma(0)}{2}\mathrm{Tr}\Bigl(\sigma^{1/2}O^\dagger\sigma^{1/2}
O\,S(0)^2\Bigr).
\end{align*}
On the other hand,
\begin{align*}
\|[S(0),O]\|^2_{2,\sigma}
&= \mathrm{Tr}\Bigl(\sigma^{1/2}\bigl(S(0)O - OS(0)\bigr)^\dagger
\sigma^{1/2}\bigl(S(0)O - OS(0)\bigr)\Bigr) \\
&= \mathrm{Tr}\Bigl(\sigma^{1/2}(O^\dagger S(0) - S(0)O^\dagger)
\sigma^{1/2}(S(0)O - OS(0))\Bigr).
\end{align*}
Using Lemma~\ref{lem:commute} to commute $S(0)$ through $\sigma^{1/2}$
and expanding the product yields
\[
\|[S(0),O]\|^2_{2,\sigma} = \mathrm{Tr}\Bigl(\sigma^{1/2}O^\dagger
\sigma^{1/2}S(0)^2 O\Bigr) - 2\,\mathrm{Tr}\Bigl(\sigma^{1/2}O^\dagger
\sigma^{1/2}S(0)\,O\,S(0)\Bigr) + \mathrm{Tr}\Bigl(\sigma^{1/2}O^\dagger
\sigma^{1/2}O\,S(0)^2\Bigr).
\]
Comparing with the expression for $\mathcal{E}^{(0)}_\sigma(O)$ gives
$\mathcal{E}^{(0)}_\sigma(O) = \frac{\gamma(0)}{2}
\|[S(0),O]\|^2_{2,\sigma}$, as claimed.

\end{document}
